% 黄道十二宫（综述）
% license CCBYSA3
% type Wiki

本文根据 CC-BY-SA 协议转载翻译自维基百科\href{https://en.wikipedia.org/wiki/Astrological_sign}{相关文章}。

\begin{figure}[ht]
\centering
\includegraphics[width=6cm]{./figures/e66f90583cecd88b.png}
\caption{出自《贝里公爵的豪华时祷书》的解剖学黄道人} \label{fig_Astsi_1}
\end{figure}
在西方占星术中，占星星座即黄道，被划分为十二个各占30度的区段，这些区段由太阳在地球天空中所观测到的360度轨道路径所穿越。星座的顺序从春季的第一天开始，即被称为“白羊宫起点”的位置，也就是春分点。占星星座依次为白羊座、金牛座、双子座、巨蟹座、狮子座、处女座、天秤座、天蝎座、人马座、摩羯座、宝瓶座和双鱼座。西方黄道起源于巴比伦占星术，后来受到希腊化文化的影响。每一个星座均以太阳每年穿越天空时经过的一个星座命名。这一观测在简化且流行的太阳星座占星术中被尤为强调。几个世纪以来，由于地球自转轴进动$^{[1]}$，西方占星术中的黄道分区逐渐与其所命名的星座发生偏离，而印度占星术的测量方法则对这种偏移进行了修正。$^{[2]}$ 占星术（即一种基于天象的预兆体系）在中国和西藏文化中也有所发展，但这些占星体系并不以黄道为基础，而是涉及整个天空。

占星术是一种伪科学。$^{[3]}$ 对其理论基础$^{[4]}$以及对相关主张的实验验证$^{[5]}$的科学研究表明，它既不具有科学有效性，也缺乏解释力。对于人格特征与出生月份之间表面相关性的现象，存在更为合理的解释，例如人类季节性出生所产生的影响。

根据占星学的观点，天体现象与人类活动之间遵循“上如其下”的原则，因此星座被认为代表着具有特征性的表达模式。$^{[6]}$ 直到19世纪之前，科学天文学与西方占星术一样，均使用相同的黄道分区。

目前，不同的占星体系采用了多种测量和划分天空的方法，尽管黄道星座的名称与符号传统在很大程度上保持一致。西方占星术以春分点和至点（即热带年中昼夜等长、最长与最短的时刻）为基准进行测量，而印度占星术则沿着赤道平面进行测量，对应的是恒星年。
\subsection{西方黄道星座}  
\subsubsection{历史}
\begin{figure}[ht]
\centering
\includegraphics[width=6cm]{./figures/c64ed924a06babf0.png}
\caption{十二个黄道星座。每一个圆点标示一个星座的起点，相邻星座之间相隔30°。天赤道与黄道的交点定义了分点：白羊宫起点（）与天秤宫起点（）。包含天极与黄极的大圆（P 与 P′）在黄道上分别与 0° 巨蟹宫（）和 0° 摩羯宫（）相交。在该示意图中，太阳被示意性地放置在宝瓶宫起点（）。} \label{fig_Astsi_2}
\end{figure}
西方占星术是希腊化占星术的直接延续，后者在公元2世纪由托勒密记录于《四书》（Tetrabiblos）之中。希腊化占星术本身又部分建立在巴比伦传统的概念之上。具体而言，将黄道划分为十二个相等区段这一思想结构源自巴比伦。$^{[7]}$ 这种黄道划分起源于古代汇编文献《MUL.APIN》中所记载的巴比伦“理想历法”，以及其与巴比伦阴历的结合，后者在《MUL.APIN》中被表述为“月亮之路”。$^{[8]}$ 从某种意义上说，黄道正是对一种理想化阴历体系的抽象化表达。

至公元前4世纪，巴比伦天文学及其天象征兆体系对古希腊文化产生了影响；而到公元前2世纪末，埃及天文学同样发挥了重要作用。这一发展不同于美索不达米亚传统，其结果是占星学的重心转向对个体出生星图的关注，并催生了本命占星术（Horoscopic astrology），引入了上升点（即出生时刻黄道升起的度数）以及十二宫的概念。将占星星座与恩培多克勒的四种古典元素相联系，也是刻画十二星座特征的另一项重要发展。

至公元2世纪时，希腊化占星传统的整体体系由托勒密在《四书》中加以系统阐述。该著作不仅是后世西方天文学传统的奠基性文本，也对印度和伊斯兰世界产生了深远影响，并在近十七个世纪中一直作为权威参考，因为后续传统几乎未对其核心教义作出实质性修改。
\subsubsection{西方占星对应表}  
下表展示了十二个占星星座的大致日期范围，以及每个星座在古典体系$^{[9]}$和现代体系$^{[10]}$中的守护星。按照定义，白羊宫始于白羊宫起点，即太阳位于三月春分点时的位置。春分的具体日期每年略有变化，但始终介于3月19日至3月21日之间。因此，白羊宫的起始日期，以及其他所有星座的起始日期，都会随年份略有变动。下述西方占星表列举了天球经度的十二个分区，并使用拉丁名称表示。经度区间在定义上，对第一个端点（a）为闭区间，对第二个端点（b）为开区间——例如，经度30°是金牛宫的起点，而不属于白羊宫。在天文学文献中，这些星座有时会以0至11的编号代替其符号来表示。
\begin{figure}[ht]
\centering
\includegraphics[width=14.25cm]{./figures/f3eef50262357bdf.png}
\caption{} \label{fig_Astsi_4}
\end{figure}
十二个星座以环形方式排列，形成一组彼此对立的关系，这些对立与不同的哲学极性属性相关联。在西方占星术中，火象与风象元素通常呈180度对立，土象与水象元素亦是如此。$^{[12]}$ 并非所有占星体系都采用四元素结构，例如《创造之书》（Sepher Yetzirah）仅描述了从一个中心神性源头发出的三种元素。$^{[13]}$ 春季星座与秋季星座相对，冬季星座与夏季星座相对，反之亦然。$^{[14][15][16][17]}$
\begin{itemize}
\item 白羊座与天秤座相对  
\item 金牛座与天蝎座相对  
\item 双子座与人马座相对  
\item 巨蟹座与摩羯座相对  
\item 狮子座与宝瓶座相对  
\item 处女座与双鱼座相对  
\end{itemize}
\textbf{极性}
在西方占星术中，极性将黄道分为两半，用以指代某一星座能量取向是积极还是消极，并由此衍生出一系列相关属性。$^{[18]}$ 积极极性的星座亦被称为主动、阳性、外放或男性星座，它们是黄道中序号为奇数的六个星座：白羊座、双子座、狮子座、天秤座、人马座和宝瓶座。积极星座构成火象与风象的三分组。$^{[19][20]}$ 消极极性的星座亦被称为被动、阴性、接纳或女性星座，$^{[19]}$ 它们是黄道中序号为偶数的六个星座：金牛座、巨蟹座、处女座、天蝎座、摩羯座和双鱼座。消极星座构成土象与水象的三分组。$^{[20]}$

\textbf{三种模式} 

某一星座的模式（modality 或 mode）指的是它在所处季节中的位置。四种元素各自以三种模式表现：本位（cardinal）、固定（fixed）和变动（mutable）。$^{[21]}$ 每一种模式涵盖四个星座，亦称为四分组（Quadruplicities）。$^{[22][23]}$ 例如，在北半球，白羊座位于春季的第一个月，因此占星学实践者将其描述为本位模式。$^{[24]}$ 元素与模式的结合赋予了每个星座独特的性格特征。例如，摩羯座是土象中的本位星座，这使其在物质世界（土象）中体现出行动力（本位模式）的关联。$^{[25][26][27]}$
\begin{figure}[ht]
\centering
\includegraphics[width=14.25cm]{./figures/9ef66be0c36c87d0.png}
\caption{} \label{fig_Astsi_5}
\end{figure}
\textbf{四元素的三分组}
\begin{figure}[ht]
\centering
\includegraphics[width=6cm]{./figures/247f42f2fabb5214.png}
\caption{2012年5月16日行星在各星座中的位置。星座按其对应的元素进行着色。每颗行星以一个符号表示，标注在其所在星座内对应的黄经位置旁。还可附加符号以表示视逆行（℞），或视为静止的瞬间（即由逆行转为顺行，或由顺行转为逆行：S）。} \label{fig_Astsi_6}
\end{figure}
希腊哲学家恩培多克勒在公元前5世纪提出火、土、气和水四种元素。他将宇宙的本性解释为两种对立原则——爱与争——之间的相互作用，这两种原则操控元素形成不同的混合，从而产生万物各异的性质。他认为所有元素彼此平等，具有相同的年代，各自统治其所属的领域，并拥有独立的个性特征。恩培多克勒指出，那些体内元素比例近乎均衡而生的人更为聪慧，并且具有最为精确的感知能力。$^{[31][32]}$

这些元素分类被称为“三分组”，因为每一种古典元素都与三个星座相关联。$^{[22][23]}$ 四种占星元素也被视为与希波克拉底的人格类型直接对应（多血质 = 气；胆汁质 = 火；忧郁质 = 土；黏液质 = 水）。一种较为现代的观点将元素视为“经验的能量物质”，$^{[33]}$ 下表试图通过关键词来概括它们的特征。$^{[34][35]}$ 随着元素概念的重要性不断提升，一些占星学家开始在解读本命盘时，首先研究行星（尤其是太阳和月亮）在各元素中的分布平衡，以及星盘中角点的位置。$^{[36]}$
\begin{figure}[ht]
\centering
\includegraphics[width=14.25cm]{./figures/f90238ee947522af.png}
\caption{} \label{fig_Astsi_7}
\end{figure}
\textbf{天体主宰}
\begin{figure}[ht]
\centering
\includegraphics[width=6cm]{./figures/fd8b2cbc306296cc.png}
\caption{1716年《学术论丛》（Acta Eruditorum）表格插图中对西方占星星座的表现} \label{fig_Astsi_8}
\end{figure}
守护关系（Rulership）指行星与其所对应的星座和宫位之间的关联。$^{[38]}$ 传统的守护关系如下所示：$^{[9][39]}$
\begin{itemize}
\item 白羊座：火星  
\item 金牛座：金星  
\item 双子座：水星  
\item 巨蟹座：月亮  
\item 狮子座：太阳  
\item 处女座：水星  
\item 天秤座：金星  
\item 天蝎座：传统上为火星，20世纪起引入冥王星  
\item 人马座：木星  
\item 摩羯座：土星  
\item 宝瓶座：传统上为土星，20世纪起引入天王星  
\item 双鱼座：传统上为木星，20世纪起引入海王星
\end{itemize}
\textbf{尊贵与失势、擢升与陷落} 

占星学中的一种传统信念被称为“本质尊贵”（essential dignity），其核心思想是：由于太阳、月亮与行星的基本性质被认为在某些星座中彼此和谐，因此它们在这些星座中的力量与效能更为强大；相反，在某些星座中，它们被认为运作乏力或受阻，因为其本性被视为相互冲突。这些范畴包括尊贵（Dignity）、失势（Detriment）、擢升（Exaltation）与陷落（Fall）。
\begin{itemize}
\item \textbf{尊贵与失势：}当一颗行星位于其所守护的星座中时，被认为其力量得到增强，或处于“尊贵”状态。换言之，它被认为在该星座中行使守护权。例如，月亮位于巨蟹座被视为“强势”（尊贵良好）。若一颗行星位于其所守护（或尊贵）星座的对宫，则被认为力量减弱，处于“失势”状态（例如月亮位于摩羯座）。这种情况有时也被称为“虚弱”。$^{[40]}$  
\end{itemize}
在传统占星术中，除守护权之外，还承认其他层级的尊贵形式，包括擢升、三分性、界限（terms 或 bounds）以及面或旬（face 或 decan）。这些层级共同用来描述行星的本质尊贵，即其真实本性的质量或能力。$^{[40]}$
\begin{itemize}
\item \textbf{擢升与陷落：}当一颗行星位于其擢升的星座中时，其力量同样会得到增强。在传统的时辰占星中，这被视为仅次于守护权的一种尊贵状态。擢升被认为赋予行星一种“受尊敬宾客”的尊位：备受关注，但在权力上有所限制。行星处于擢升状态的例子包括：土星（天秤座）、太阳（白羊座）、金星（双鱼座）、月亮（金牛座）、水星（处女座，尽管对此存在不同意见）、火星（摩羯座）、木星（巨蟹座）。若行星位于其擢升星座的对宫，则处于“陷落”状态，其力量被削弱，甚至可能比失势更为严重。$^{[40]}$ 对于两颗外行星（超土星行星）应被视为在哪些星座中擢升，学界仍存在分歧。$^{[41]}$
\end{itemize}
\begin{figure}[ht]
\centering
\includegraphics[width=14.25cm]{./figures/3b0126e155258735.png}
\caption{} \label{fig_Astsi_9}
\end{figure}
除本质尊贵之外，传统占星学家还会考虑行星的偶然尊贵（Accidental dignity）。这指的是行星在所考察星盘中所处宫位的位置。偶然尊贵体现的是行星“发挥作用的能力”。例如，月亮位于巨蟹座时，因守护权而具有尊贵，但若其落在第十二宫，则几乎没有空间来表达其良好的本性。$^{[42]}$ 第十二宫属于衰退宫（cadent house），第三宫、第六宫和第九宫亦是如此，行星位于这些宫位通常被认为力量薄弱或受损。相反，月亮若位于第一、第四、第七或第十宫，则更有能力发挥作用，因为这些属于角宫（angular houses）。行星位于继起宫（succedent houses，即第二、第五、第八和第十一宫）时，一般被认为其作用能力处于中等水平。除偶然尊贵之外，还存在一系列偶然失势（Accidental Debilities），例如逆行、日光下（Under the Sun's Beams）、灼烧（Combust）等情况。

\textbf{附加分类}
\begin{figure}[ht]
\centering
\includegraphics[width=10cm]{./figures/bab047d44858754b.png}
\caption{现代星座的赤纬—赤经等距圆柱投影图，虚线表示黄道。星座按其所属家族及确立年份进行颜色编码。（详细视图）} \label{fig_Astsi_10}
\end{figure}
每一个星座都可以被划分为三个各占 10° 的区段，称为旬（decans 或 decanates），不过这一划分方式如今已较少使用。第一个旬被认为最能体现该星座自身的本质，并由该星座的守护行星所统治。$^{[43]}$ 第二个旬由同一三分性中下一个星座的守护行星作为次级守护。最后一个旬则由同一三分性中再下一个星座的守护行星作为次级守护。$^{[44]}$

虽然星座的元素与模式已经足以界定其特性，但它们还可以进一步分组以指示其象征意义。前四个星座——白羊座、金牛座、双子座和巨蟹座——构成“个人星座”组。接下来的四个星座——狮子座、处女座、天秤座和天蝎座——构成“人际星座”组。黄道最后的四个星座——人马座、摩羯座、宝瓶座和双鱼座——构成“超个人星座”组。$^{[45]}$

丹恩·鲁德亚（Dane Rudhyar）提出了热带黄道的基本因素，$^{[46]}$ 并被用于 RASA 占星学院的课程体系之中。热带黄道以季节性因素为基础，与恒星黄道（以星座为基础）相对。其主要季节因素基于一年之中日照与黑暗比例的变化。第一个因素是所选时间是否位于白昼增加的一年半周期，或位于黑暗增加的一年半周期。第二个因素是所选时间是否位于白昼多于黑暗的半年，或黑暗多于白昼的半年。第三个因素则是所选时间落在四季中的哪一季，而四季正是由前两个因素共同定义的。因此$^{[47][48]}$
\begin{itemize}
\item 春季是白昼增加且白昼多于黑暗的时期。  
\item 夏季是黑暗增加且白昼多于黑暗的时期。  
\item 秋季是黑暗增加且黑暗多于白昼的时期。  
\item 冬季是白昼增加且黑暗多于白昼的时期。
\end{itemize}
\subsubsection{西方星座图库}
\begin{figure}[ht]
\centering
\includegraphics[width=14.25cm]{./figures/d46708c77f7f01c9.png}
\caption{} \label{fig_Astsi_12}
\end{figure}
\subsection{印度占星术}
在印度占星术（Jyotiṣa）中，其宇宙论框架基于“五大元素”（Pancha Mahābhūta）：火（Agni）、土（Pṛthvī）、气（Vāyu）、水（Jala）与以太（Ākāśa）。火的主宰为火星，水星对应土地（土），土星对应空气，金星对应水，木星对应以太。

五颗可见行星（不包括太阳与月亮）各自与其中一种元素相关联，并作为其统治原则：
\begin{itemize}
\item 火星（Maṅgala）为火（Agni）之主  
\item 水星（Budha）主宰土（Pṛthvī）  
\item 土星（Śani）主宰气（Vāyu）  
\item 金星（Śukra）主宰水（Jala）  
\item 木星（Bṛhaspati）主宰以太（Ākāśa）$^{[49]}$
\end{itemize}
印度占星术承认十二个黄道星座（Rāśi），$^{[50]}$ 与西方占星术中的星座相对应。两种体系中，星座与元素之间的对应关系是相同的。
\subsubsection{宿}  
宿（天城文：नक्षत्र，梵语 nakshatra）又称月宿，是印度天文学与占星术（Jyotiṣa）中依据显著恒星划分天空的27个区域之一。$^{[51]}$ “Nakshatra”一词在梵语、卡纳达语、图卢语、泰米尔语以及俗语（Prakrit）中，均直接指恒星本身。
\subsection{中国黄道}  
中国占星符号以年、农历月份以及一天中两个小时的时段（亦称时辰）为循环单位运作。中国黄道的一大特征在于其与中国占星术中的五行（木、火、金、水、土）相结合，形成一个60年的循环周期。$^{[52]}$ 尽管如此，一些研究指出，中国的十二年周期与黄道星座之间存在明显联系：周期中的每一年都对应木星的某种位置分布。例如，在蛇年，木星位于双子座；在马年，木星位于巨蟹座，依此类推。因此，中国的十二年历法可被视为一种太阳—月亮—木星历法。
\subsubsection{黄道象征}
下表显示了十二星座及其属性。
\begin{figure}[ht]
\centering
\includegraphics[width=8cm]{./figures/3aca0f8778275f52.png}
\caption{中国生肖占星术中的符号} \label{fig_Astsi_13}
\end{figure}
\subsubsection{十二星座}
\begin{figure}[ht]
\centering
\includegraphics[width=6cm]{./figures/c891c8bb599317ef.png}
\caption{显示24个基本方位及其对应星座符号的图表。} \label{fig_Astsi_14}
\end{figure}
在中国占星术中，由十二种动物组成的黄道代表了十二种不同的人格类型。黄道传统上以鼠为首，关于中国黄道起源的诸多传说解释了这一排序的缘由。当十二生肖与四种元素结合，纳入六十年历法体系时，它们传统上被称为十二地支。中国黄道遵循阴阳合历的中国历法，$^{[53]}$ 因此每个月中“交替”的日期（即一个生肖转换为下一个生肖的时间）每年都会有所不同。以下依次列出十二生肖。$^{[54]}$
\begin{enumerate}
\item 子　鼠（阳性，第1三合，固定元素：水）：  
鼠年包括1900、1912、1924、1936、1948、1960、1972、1984、1996、2008、2020、2032年。鼠亦对应一年中的特定月份。鼠时为晚上11点至凌晨1点。
\item 丑　牛（阴性，第2三合，固定元素：土$^{[55]}$）：  
牛年包括1901、1913、1925、1937、1949、1961、1973、1985、1997、2009、2021、2033年。牛亦对应一年中的特定月份。牛时为凌晨1点至3点。
\item 寅　虎（阳性，第3三合，固定元素：木）：  
虎年包括1902、1914、1926、1938、1950、1962、1974、1986、1998、2010、2022、2034年。虎亦对应一年中的特定月份。虎时为凌晨3点至5点。
\item 卯　兔（阴性，第4三合，固定元素：木）：  
兔年包括1903、1915、1927、1939、1951、1963、1975、1987、1999、2011、2023、2035年。兔亦对应一年中的特定月份。兔时为凌晨5点至7点。
\item 辰　龙（阳性，第1三合，固定元素：土$^{[55]}$）：  
龙年包括1904、1916、1928、1940、1952、1964、1976、1988、2000、2012、2024、2036年。龙亦对应一年中的特定月份。龙时为上午7点至9点。
\item 巳　蛇（阴性，第2三合，固定元素：火）：  
蛇年包括1905、1917、1929、1941、1953、1965、1977、1989、2001、2013、2025、2037年。蛇亦对应一年中的特定月份。蛇时为上午9点至11点。
\item 午　马（阳性，第3三合，固定元素：火）：  
马年包括1906、1918、1930、1942、1954、1966、1978、1990、2002、2014、2026、2038年。马亦对应一年中的特定月份。马时为上午11点至下午1点。
\item 未　羊（阴性，第4三合，固定元素：土$^{[55]}$）：  
羊年包括1907、1919、1931、1943、1955、1967、1979、1991、2003、2015、2027、2039年。羊亦对应一年中的特定月份。羊时为下午1点至3点。
\item 申　猴（阳性，第1三合，固定元素：金）：  
猴年包括1908、1920、1932、1944、1956、1968、1980、1992、2004、2016、2028、2040年。猴亦对应一年中的特定月份。猴时为下午3点至5点。
\item 酉　鸡（阴性，第2三合，固定元素：金）：  
鸡年包括1909、1921、1933、1945、1957、1969、1981、1993、2005、2017、2029、2041年。鸡亦对应一年中的特定月份。鸡时为下午5点至7点。
\item 戌　狗（阳性，第3三合，固定元素：土$^{[55]}$）：  
狗年包括1910、1922、1934、1946、1958、1970、1982、1994、2006、2018、2030、2042年。狗亦对应一年中的特定月份。狗时为晚上7点至9点。
\item 亥　猪（阴性，第4三合，固定元素：水）：  
猪年包括1911、1923、1935、1947、1959、1971、1983、1995、2007、2019、2031、2043年。猪亦对应一年中的特定月份。猪时为晚上9点至11点。
\end{enumerate}
\subsubsection{五行}  

\begin{itemize}
\item 木：木性之人道德感强，自信，性格外向且善于合作，兴趣广泛多样，目标理想主义。木所对应的方位为东方，季节为春季，因此它是虎与兔这两个生肖的固定元素。$^{[55]}$
\item 火：火性之人具备领导才能，充满活力与激情，果断、自信、积极且富有进取性。火所对应的方位为南方，季节为夏季，因此它是蛇与马这两个生肖的固定元素。$^{[55]}$
\item 土：土性之人严肃、理性且有条理，聪慧、客观，善于规划。土所对应的方位为中央，其季节为四季交替之时。它是牛、龙、羊与狗这四个生肖的固定元素。$^{[55]}$
\item 金：金性之人真诚，价值观与立场坚定，意志坚强，且言辞雄辩。金所对应的方位为西方，季节为秋季。它是猴与鸡这两个生肖的固定元素。$^{[55]}$
\item 水：水性之人富有说服力，直觉敏锐，富于同理心。水性之人客观，且常因善于提供建议而受到他人求助。水所对应的方位为北方，季节为冬季。它是鼠与猪这两个生肖的固定元素。$^{[55]}$
\end{itemize}
五行与十二生肖共同构成一个六十年的历法循环。五行在历法中以阴与阳两种形式出现，被称为十天干。公历中体现的阴阳划分意味着，年份数字以偶数结尾者为阳年（象征阳刚、主动与光明），以奇数结尾者为阴年（象征阴柔、被动与黑暗），这一判断以当年中国新年已经到来为前提。$^{[55]}$
\subsection{参见}  
\begin{itemize}
\item 季节性出生对人类的影响  
\item 中国生肖  
\item 占星术术语表
\end{itemize}
\subsection{注释}
 博布里克（2005），第10、23页。
 约翰森（2004）。
 斯文·奥韦·汉松；爱德华·N·扎尔塔。“科学与伪科学”。斯坦福哲学百科全书。检索日期：2012年7月6日例如，人们普遍认同，创世论、占星术、顺势疗法、基尔利安摄影、灵摆学、不明飞行物研究、古代宇航员理论、否认大屠杀、维利科夫斯基式灾变论以及气候变化否认主义都属于伪科学。
 维什韦沙瓦拉（1989）。S.K. 比斯瓦斯；D.C.V. 马利克；C.V. 维什韦沙瓦拉（编）。宇宙视角：献给M.K.V. 巴普纪念的论文集（第1版）。英国剑桥：剑桥大学出版社。ISBN 978-0-521-34354-1.
 卡尔森，肖恩（1985）。“占星术的双盲测试” （PDF）. 自然. 318 (6045): 419–425. 文献标识码:1985Natur.318..419C. DOI:10.1038/318419a0. S2CID 5135208. 已归档 （PDF）于2019年2月16日自原始网页存档。已获取2021-05-26.
 马约（1979），第35页。
 萨克斯（1948），第289页。苏美尔文献中对天体“征兆”的零星提及不足以证明存在苏美尔黄道十二宫，参见罗克伯格（1998），第ix页。
 霍夫曼，苏珊娜·M.（2017）。希帕尔科斯的天球仪：巴比伦与希腊天文学之间的纽带？维斯巴登：施普林格·施佩克特姆。ISBN 978-3-658-18683-8. OCLC 992119256.
 卡罗尔·威尔斯（2007）。“统治权”。占星学现在。于2020-10-28从原文存档。检索日期：2007-11-29
 霍恩（1978），第21页。
 杰里米·B·塔特姆，《黄道十二宫的星座与星群》存档于2023年6月1日，见互联网档案馆》，加拿大皇家学会期刊,104（2010），103—104页。
 “反对派”.dichesegnosei.it（意大利语）。已于2017年7月16日从原始网页存档。检索日期：2017年7月27日
 “WoC :: 塞弗·耶齐拉”. www.workofthechariot.com. 已存档自原始页面于2023-01-20。检索日期2023-01-20.
 “空气元素”.dichesegnosei.it（意大利语）。已于2017-07-29从原始网页存档。检索日期：2017-07-27
 “火之元素”.dichesegnosei.it（意大利语）。已于2017-08-08存档自原始网页。检索日期：2017-07-27
 “水元素”.dichesegnosei.it（意大利语）。已于2017年7月22日从原始网页存档。检索日期：2017年7月27日
 “地球元素”.dichesegnosei.it（意大利语）。已于2017年6月30日从原始网页存档。检索日期：2017年7月27日
 霍尔，朱迪（2005）。《占星术圣经：黄道十二宫终极指南》。斯特林出版公司，Inc. 第137页。ISBN 978-1-4027-2759-7.
 斯坦登，安东尼（1975）。“人格是否存在星象效应”。《心理学杂志》89（2）：259–260。DOI：10.1080/00223980.1975.9915759PMID：1151896已于2020年7月25日存档自原始网页。检索于2011年11月20日
 范·罗伊，扬·J·F.（1993）。“内向—外向：占星术与心理学的比较”。荷兰莱顿大学心理学系. 16 (6): 985–988. doi:10.1016/0191-8869(94)90243-7.
 阿罗约（1989），第29页。
 佩利蒂埃，罗伯特，和莱昂纳德·卡塔尔多（1984）。做自己的占星师。伦敦：潘图书。第43–44页
 波滕格，玛丽莎（1991）。天文学基础。圣地亚哥：ACS出版公司。第31–36页。
 伍德韦尔2019, 第26页。
 霍恩（1978），第75页
 蒂尔尼，比尔（2001）。十二宫全览：探索占星术的十二星座。莱韦林全球出版社。ISBN 978-0-7387-0111-0.
 伍德韦尔，唐娜（2019-11-19）。占星术词典：从A到Z的宇宙知识。西蒙与舒斯特出版社。ISBN 978-1-5072-1144-1.
 如塞法里尔《占星术手册》巴西版（1988年）所用，由里约热内卢新前沿出版社S/A出版
 霍恩（1978），第40页
 阿罗约（1989），第30页
 查理·希金斯（1997）。《占星术与四大元素》。于2010-01-03从原文存档。于2009-12-27获取。
 阿杰·莫汉（2002）。“占星术与四大元素”。于2020-05-14从原文存档。检索日期：2017-08-03.
 阿罗约（1989），第27页。
 阿罗约（1989），第30至34页
 霍恩（1978），第42页
 阿罗约（1975）第131至140页。
 来自炼金术象征符号的符号。
 霍恩（1978），第22页
 沙莫斯，杰弗里（2013）。“占星术作为一种社会框架：‘行星之子’，1400–1600年”。宗教、自然与文化研究期刊. 7 (4): 434–460. doi:10.1558/jsrnc.v7i4.434.
 “占星术术语表”。Logos，Asaa 1998–2004。2007年11月26日。“Logos专业占星学——占星术术语表”。于2008年1月22日存档自原文。检索日期：2007年11月29日。
 霍恩（1978），第144页
 “意外的尊严”.www.gotohoroscope.com。占星学词典 1998–2007。2007年11月26日。已存档自原始网页于2008年1月3日保存2007年11月29日检索到
 霍恩（1978），第87页
 霍恩（1978），第88页
 “占星术导论。”Spiritsingles.com，2007年11月25日。“占星运势图”。于2007-12-09从原文存档。检索日期：2007-11-29。
 鲁迪亚尔（1943）
 阿姆斯特朗，罗宾（2009）。“星座与宫位”。RASA占星学院。于2012-03-20从原始网页存档。检索日期：2011-02-20。
 阿姆斯特朗，罗宾（2009）。《易经：变化的序列》。RA出版社。于2011年7月24日从原始版本存档。检索日期：2011年2月20日。
 “根据占星术，能量及其相关元素是什么？”. www.melooha.com。检索日期：2025-09-21.
 萨顿（1999）第74至92页。
 萨顿（1999），第168页。
 “《年鉴》‘农历’黄道十二宫以立春为界存在错位？——中国网”，2009年2月16日。于2011年6月14日存档自原始网页。检索日期：2011年1月5日。
 诺瓦克，萨拉。“中国十二生肖”。面孔. 34: 42.
 西奥多拉·劳,同上, 第2–8页，第30–5页，第60–4页，第88–94页，第118–24页，第148–53页，第178–84页，第208–13页，第238–44页，第270–78页，第306–12页，第338–44页，2005年
 中国占星术：探索东方黄道十二宫，作者：谢莉·吴